(1) 球形新核增加面积 \(A=4\pi r^2\)，耗费表面自由能 \(4\pi\gamma r^2\)；由于内部转化为热力学上更优的新相，体积自由能降低 \((4\pi/3)\Delta g\,r^3\)。因此 \(\Delta G(r)=4\pi\gamma r^2-(4\pi/3)\Delta g\,r^3\)。若缺少正的表面项，任何无穷小核都会生长，与实际存在的过饱和亚稳态不符。
(2) 求导 \(\Delta G'(r)=8\pi\gamma r-4\pi\Delta g\,r^2=4\pi r(2\gamma-\Delta g\,r)\)，非零驻点为 \(r_*=2\gamma/\Delta g\)。二阶导在该点为 \(\Delta G''(r_*)=-8\pi\gamma<0\)，故临界核是自由能障碍的最大点而非稳定局部最小。代回得到 \(\Delta G_*=16\pi\gamma^3/(3\Delta g^2)\)，单位为能量。
(3) 当 \(\Delta g\to0^+\) 时 \(r_*\propto\Delta g^{-1}\)、\(\Delta G_*\propto\Delta g^{-2}\)，弱过饱和使成核异常稀少。在经典 Arrhenius–Kramers 近似下成核速率 \(J\simeq J_0\exp[-\Delta G_*/(k_BT)]\)，\(J_0\) 取决于分子输运、附着动力学和位点数。该公式描述的是越垒稀有事件，不能由静态自由能曲线单独确定前因子。
(4) 异相成核在可润湿基底上往往只需形成球冠，界面几何减少能垒；理想平滑基底的形状因子为 \(f(\vartheta)=(2+\cos\vartheta)(1-\cos\vartheta)^2/4\)，有 \(\Delta G_*^{\rm het}=f(\vartheta)\Delta G_*^{\rm hom}\)。当临界核只有几个分子、\(\gamma\) 随曲率变化显著、界面各向异性或存在弹性应变时，连续球面近似及常数 \(\gamma\) 不再可靠。
